\chapter{Order, optimisation and simulation}
\label{ch:order}

A laser job is a list of moves, and the order of those moves decides how long the
job takes, how cleanly it cuts, and whether small parts stay where you put them.
This chapter is about the planner --- the part of MeerK40t that turns your
operations into instructions --- and about the two windows that let you inspect
the result before it reaches the machine.

\section{The plan}

Between pressing \btn{Start} and the machine moving, MeerK40t runs a fixed
sequence of stages. You have already met them in the console output; here is what
each one does:

\begin{center}\small
\begin{tabular}{@{}L{2.3cm}L{11.9cm}@{}}
\toprule
\textbf{Stage} & \textbf{What happens} \\
\midrule
\emph{copy} & The operations are copied into a plan; references to elements are
replaced by real copies, so the plan is immune to later edits. \\
\emph{preprocess} & Scene coordinates are converted to device coordinates, and
validation operations are inserted. \\
\emph{validate} & The validation operations run --- checking that what follows can
actually be executed. \\
\emph{blob} & Operations and elements become \emph{cutcode}: lines, arcs, raster
strokes and moves. This is where an operation's settings finally become machine
instructions. \\
\emph{preopt} & Optimisation commands are injected --- inner-first ordering,
travel reduction, sequencing. \\
\emph{optimize} & Those commands run, reordering the cutcode. \\
\emph{spool} & The finished job is wrapped up and handed to the spooler, which
feeds the driver. \\
\bottomrule
\end{tabular}
\end{center}

\noindent You can drive the same stages by hand. The \btn{Execute Job} window
(\menu{Tools>Execute Job}) is a staged wizard whose button changes label as you
advance: \btn{Copy}, \btn{Preprocess Operations}, \btn{Validate}, \btn{Convert
data}, \btn{Preprocess Optimizations}, \btn{Optimize}, \btn{Spool}. The
\tabname{Plan} tab of the \panel{Laser-Control} pane shows the plan itself ---
columns \emph{\#}, \emph{Plan}, \emph{Status} and \emph{Content} --- with
\btn{Update}, \btn{Export} and \btn{Spool}, and a \btn{Hold} checkbox which
preserves the plan between runs so you can inspect and rerun it. Right-clicking a
plan offers \btn{Clear}, \btn{Update}, \btn{Update (background)}, \btn{Export},
\btn{Spool} and \btn{Details}.

\begin{behindbox}[Why you would ever do it by hand]
Two reasons. First, diagnosing: if a job behaves oddly, stopping after
\emph{blob} shows you the cutcode the plan produced, before optimisation had a
chance to hide the problem. Second, iterating: with \btn{Hold} on, a plan can be
optimised repeatedly and compared in the simulation window without rebuilding
from scratch each time.
\end{behindbox}

\section{The optimisation settings}

\menu{Edit>Settings>Preferences}, or the \tabname{Optimize} tab of the laser
panel, exposes the planner's settings:

\begin{center}\small
\begin{tabular}{@{}L{4.3cm}L{1.7cm}L{8.2cm}@{}}
\toprule
\textbf{Label} & \textbf{Default} & \textbf{Effect} \\
\midrule
\btn{Reduce Travel Time} & on & Reorders burns to shorten the distance travelled
with the laser off. The single biggest time saver on most jobs. \\
\btn{Burn Inner First} & on & Cuts inner contours before outer ones, so a part is
still held by the surrounding material when its own holes are being cut. \\
\btn{Group Inner Burns} & off & Completes one part --- inner and outer --- before
starting the next, instead of doing all inner cuts first. Slower, but safer on
sheets where parts could shift. \\
\btn{Merge Passes} & off & Optimises the passes of an operation together rather
than pass by pass. \\
\btn{Merge Operations} & off & Optimises across operations instead of within
each. \\
\btn{Burn Complete Subpaths} & on & Finishes each open subpath in one go. \\
\btn{Combine path segments} & off & Stitches nearly-touching path ends so they
burn as one continuous move. With \emph{Tolerance} (default \cmd{0}). \\
\btn{Cluster raster objects} & on & Groups rasters that do not overlap into
separate optimisation groups, with a \emph{Margin} (default \cmd{1mm}). \\
\btn{Keep effect lines together} & on & Treats hatch lines as one shape, so they
are burned together rather than visited one by one. \\
\btn{Optimize internally} & off & Optimises the lines inside a hatch
individually. \\
\btn{Reduce polyline details} & off & Thins redundant points, with a \emph{Level}
(\emph{Minimal} \dots\ \emph{Coarse}). A fix for paths with absurd numbers of
nodes. \\
\bottomrule
\end{tabular}
\end{center}

\noindent Two further settings have no user interface but exist and can be set
from the console: \setting{opt_nearest_neighbor} (on --- the travel algorithm
used by \emph{Reduce Travel Time}), \setting{opt_2opt} (off --- a slower
two-opt refinement), \setting{opt_start_from_position} (off --- begin from the
head's current position) and \setting{opt_closed_distance} (15 --- the distance
at which path ends are considered closed). Which strategy runs depends on the
combination chosen: with \emph{Burn Inner First} on, the inner-first optimiser
handles everything; otherwise travel reduction runs if it is enabled, and a basic
sequencing pass is the fallback.

\begin{figure}[htbp]
\centering
\begin{tikzpicture}[scale=0.9]
\draw[mkgrey!70,line width=0.8pt] (0,0) rectangle (4.2,4.2);
\draw[mkblue,line width=0.8pt] (0.5,0.5) rectangle (3.7,3.7);
\draw[mkred,line width=0.8pt] (1.1,1.1) rectangle (3.1,3.1);
\draw[mkgreen,line width=0.8pt] (1.7,1.7) rectangle (2.5,2.5);
\node[mkdark,font=\sffamily\scriptsize,anchor=west] at (4.5,3.0)
  {cut order matters, not nesting:};
\node[mkred,font=\sffamily\scriptsize,anchor=west] at (4.5,2.6)
  {1\quad innermost first};
\node[mkgreen,font=\sffamily\scriptsize,anchor=west] at (4.5,2.2)
  {2\quad then outwards};
\node[mkblue,font=\sffamily\scriptsize,anchor=west] at (4.5,1.8)
  {3\quad last};
\node[mkgrey,font=\sffamily\scriptsize,anchor=west] at (4.5,1.4)
  {4\quad outside, last of all};
\end{tikzpicture}
\caption{Why inner-first is the default. If the outer contour is cut first, the
part is held by nothing but friction while its interior is being cut --- and on
thin material, the vibration is enough to move it. Burning from the inside out
keeps the surrounding material as a clamp until the very last cut.}
\label{fig:innerfirst}
\end{figure}

\begin{tipbox}[A sensible default set for most sheet work]
\btn{Reduce Travel Time} on, \btn{Burn Inner First} on, \btn{Group Inner Burns}
off, \btn{Burn Complete Subpaths} on, \btn{Keep effect lines together} on. Build
the plan, simulate it, then change one setting and compare. Optimisation is one of
the few places where measurement beats intuition, because the numbers are right
there in the simulation window.
\end{tipbox}

\section{The Simulation window}

\btn{Simulate} on the Laser tab (right-clicking it simulates without
optimisation) opens the \panel{Simulation} window: a replay of the job, drawn on
the scene.

\begin{center}\small
\begin{tabular}{@{}L{4.0cm}L{10.2cm}@{}}
\toprule
\textbf{Control} & \textbf{What it does} \\
\midrule
Play button & Start or stop the replay. \\
\emph{Playback Speed} & Replay speed, with a percentage readout. \\
\emph{Mode} & Whether the replay advances in \emph{Steps}, \emph{Time (sec.)} or
\emph{Time (min)} --- so a four-hour job is watchable. \\
\btn{Send to Laser} & Spools the currently displayed plan. This is the
shortcut from ``looks right'' to ``running''. \\
\btn{Recalculate} / \btn{Stop calculating} & Rebuild the plan, and stop a long
rebuild. During the rebuild the button reads \emph{Calculating} and the window
says the plan is still being processed. \\
\btn{Optimize} & Whether the rebuild includes the optimisation stages. \\
\bottomrule
\end{tabular}
\end{center}

\noindent The estimates are the most useful numbers the program produces:

\begin{itemize}
  \item \emph{Laser Time} and \emph{Laser Distance} --- beam on.
  \item \emph{Travel Time} and \emph{Travel Distance} --- beam off.
  \item \emph{Total Time} and \emph{Total Distance} --- the sum.
  \item \emph{Extra Time} --- allowances for acceleration, delays and the machine's
        own overheads, and usually the reason reality is longer than the
        simulation.
\end{itemize}

\noindent Right-click on the simulation scene for the display options:
\emph{Delete cuts before}, \emph{Delete cuts after} (trim the plan to a point in
the replay), \emph{Show Background}, the three grid toggles, \emph{Show physical
dimensions}, \emph{Show travel path} (draw the jump moves --- the fastest way to
see whether optimisation did anything sensible), \emph{Raster as Image},
\emph{Simulate laser width}, and the zoom commands. The embedded operations list
in the window shows the plan's operations; the \emph{D} checkbox
\emph{decompiles} the cutplan --- making the operations visible and editable
again --- and the text field beside it lets you modify an operation parameter and
press \key{Enter} to see the effect on the simulation.

\begin{tipbox}[Watch the travel path, not the burn]
The laser path is what you expect to see; the travel path is where the time goes.
Turn on \emph{Show travel path} and look for long diagonal jumps between things
that are physically next to each other. Each one is a missing optimisation --- or
artwork that needs stitching (\cmd{stitch 0.1mm}).
\end{tipbox}

\section{The spooler and running many jobs}

\panel{Spooler} (\menu{Tools>Spooler} or the ribbon's \btn{Spooler} button; the
console form is \cmd{window toggle JobSpooler}, or \cmd{pane toggle spooler} to
work on the pane itself) is a docked pane, not a window of its own, and it is the
queue between planning and the machine.

\begin{itemize}
  \item The top row chooses the device and offers \btn{Pause}/\btn{Resume},
        \btn{Abort} and a \emph{Silent} checkbox.
  \item The queue lists each job with its \emph{\#}, \emph{Device}, \emph{Name},
        \emph{Items}, \emph{Status}, \emph{Type}, \emph{Steps}, \emph{Passes},
        \emph{Priority}, \emph{Runtime} and \emph{Estimate}. Right-click a job for
        \btn{Stop} (while running) or \btn{Remove}, \btn{Disable}/\btn{Enable},
        \emph{Finish after this loop} and \emph{Clear All}.
  \item Below the queue, \emph{Completed jobs:} is the history, with columns for
        device, name, start, end, runtime, estimate, steps, passes, status and job
        information, and a \btn{Clear History} button. Right-clicking the history
        offers \emph{Delete this entry}, \emph{Delete\dots}, \emph{All entries},
        \emph{Older than \emph{n} week}, \emph{All incomplete jobs},
        \emph{Clear history on startup}, \emph{Ignore helper jobs} and
        \emph{Write CSV} --- which writes \file{history.csv} into the program's
        working directory, and is how you keep a machine log without typing.
\end{itemize}

\noindent The \btn{Launch Spooler on Job Start} preference (on by default) brings
the spooler up automatically whenever a job starts, which is a good default: the
queue is where you see what is about to happen.

\begin{machinebox}{Judge the order by the sound, then by the result}
The simulation tells you what the planner thinks; the machine tells you the truth.
Listen for the head crossing the sheet repeatedly, and look at the first finished
part: if small pieces have shifted, you cut outer contours too early, and
\btn{Group Inner Burns} or a reordered operation list is the fix. If the edges are
charred on the last cut only, the part moved because it was no longer supported
--- a tab (\btn{Tab Edit}) solves that more reliably than any optimiser.
\end{machinebox}

\begin{tryit}{Optimise a real job, with numbers}
\begin{enumerate}[label=\arabic*.,leftmargin=1.4em]
  \item Load a job with a mixture of cutting and engraving, or build one from a
        duplicated part.
  \item Simulate it with optimisation \emph{off}. Write down \emph{Total Time},
        \emph{Travel Time} and \emph{Total Distance}.
  \item Turn optimisation on, press \btn{Recalculate}, and write the numbers down
        again. The difference is what the optimiser earned you.
  \item Turn on \emph{Show travel path} and look for the remaining long jumps.
  \item In the console, run \cmd{stitch 0.1mm} on the artwork and recalculate.
        Note the change in total distance.
  \item Try \btn{Group Inner Burns} on and compare both the numbers and a physical
        test cut. Decide for your material which one holds the parts better.
  \item Keep the winning combination as your default in the preferences, and
        record the numbers with your material notes.
\end{enumerate}
\end{tryit}
